% Source: https://artofproblemsolving.com/community/c3h2274398_geometry_notation
% Stable identity: aops-topic-2274398-q01
\textbf{(a)} For $k>0$, the image $A'$ lies on the ray from $P$ through $A$, at distance $PA'=k\cdot PA$. In vector form, $A'=P+k(A-P)$. (The centre $P$ itself is fixed.) Saying only ``$A'$ is twice as far from $P$ as $A$ is'' is not enough: all points of the circle with centre $P$ and radius $2\cdot PA$ have that property. The direction matters too: $A'$ is on the same ray from $P$ as $A$. (For $k<0$ the image is on the opposite ray, at distance $|k|\cdot PA$.)

\textbf{(b)} Using $A'=P+2(A-P)$:
\[A'=(1,2)+2(3,-3)=\boxed{(7,-4)},\qquad B'=(1,2)+2(2,3)=\boxed{(5,8)}.\]

\begin{asy}
usepackage("amsmath");
size(5.5cm);
defaultpen(fontsize(10pt));
pair P=(1,2),A=(4,-1),A2=(7,-4),B=(3,5),B2=(5,8);
draw((-0.5,0)--(8,0),gray,Arrow(4)); draw((0,-5)--(0,9),gray,Arrow(4));
draw(P--A2,dashed); draw(P--B2,dashed);
draw(A--B,linewidth(0.8)); draw(A2--B2,blue+linewidth(0.8));
dot(P);dot(A);dot(B);dot(A2,blue);dot(B2,blue);
label("$P(1,2)$",P,W); label("$A(4,-1)$",A,SW); label("$B(3,5)$",B,W);
label("$A'(7,-4)$",A2,E,blue); label("$B'(5,8)$",B2,E,blue);
label("$x$",(8,0),S); label("$y$",(0,9),W);
\end{asy}

\textbf{(c)} $B'-A'=\left(P+2(B-P)\right)-\left(P+2(A-P)\right)=2(B-A)$. Here $B-A=(-1,6)$ and $B'-A'=(-2,12)=2(-1,6)$. So the segment $A'B'$ has the same direction as $AB$, hence is parallel to it, and $A'B'=2\cdot AB=2\sqrt{37}$.
